
Permutation-equivariant architectures for weight-space learning---including Deep Weight Space (DWS) and Neural Functional Transformers (NFT)---encode different inductive biases, yet no systematic comparison has evaluated how these differences translate to task performance. This work provides the first quantitative measurement of inductive bias differences between DWS and NFT: DWS produces weight updates with coefficient of variation (CoV) 1.44 across layers, while NFT produces more uniform updates with CoV 1.35. On a synthetic accuracy prediction task, NFT achieves RMSE 82.9 compared to DWS RMSE 94.5, a 12.3\% improvement consistent with the hypothesis that global attention benefits holistic property aggregation. However, a synthetic backdoor detection task yielded chance-level performance (AUC ~0.48) for all architectures, leaving the hypothesized DWS locality advantage on local anomaly detection unverified. The architecture-task interaction effect is statistically significant (F=45616, p<0.001), but only one direction of the predicted crossover pattern was confirmed. These findings provide partial guidance for practitioners: global attention architectures are preferable for aggregate statistics tasks, while the case for locality-preserving architectures on local pattern detection remains open pending evaluation on real benchmarks.
